\section{Construction validation on new graph samples}
\label{sec:certificate-validation}
The original constructor was developed after inspecting the unresolved
instances. To test its behavior beyond those instances, we fixed a
separate protocol before executing the following study. It uses seeds
100--104 at each of 20, 40 and 60 vertices for random trees, BA graphs
with $m=2$, and ER graphs with $p=0.15,0.30$: 60 random graphs in total.
Nine structural stress cases are added: grids $P_3\square P_7$,
$P_4\square P_9$, $P_5\square P_{11}$; cycles $C_{20},C_{40},C_{60}$;
and $K_{10,10},K_{20,20},K_{30,30}$. Each graph uses $(1,1),(2,1),(3,2)$,
giving 207 instances. Embedded edge sets, source hashes and configuration
were saved before execution. No labeled graph fingerprint overlaps the
original cohort; graph isomorphism was not used as an exclusion test.
The protocol is locally fixed, not externally preregistered. All samples
and unsuccessful searches are retained, without retuning after observing results.

\paragraph{Controls and endpoints.}
Five construction methods are compared. Greedy assigns the smallest
allowed label in vertex-index order. Matching-$h$ gives the frozen
pair/singleton classes labels $0,h,2h,\ldots$. Random order and insertion
order use that same partition and seek compatible orders, with at most
100 attempts and algorithm seed 0; successful orders use spacing $k$.
For $h=k$, both ordering arms use the canonical order without search.
The final baseline applies DSATUR to $G^2$ and multiplies its colors by
$h$; it uses NetworkX's saturation-largest-first greedy implementation
of the classical coloring heuristic \cite{brelaz1979coloring}.
Thus every baseline is an explicit labeling construction, rather than a
claim of exact coloring or a reproduction of a specialized heuristic.

The shared lower bound is the maximum of the degree bound, a checked
square-clique bound, and the matching certificate described below when
applicable. Cliques are obtained by a fixed greedy search from every
starting vertex; maximum size is not assumed. Family-specific exact
formulas are not added to this common certificate pool. A positive gap
therefore means that this pool does not establish optimality, not that
the labeling is suboptimal or that its optimum is unknown in the literature.
Every proposed labeling is checked directly against the embedded graph.
Attempt counts are not equal wall-time budgets, and this experiment makes
no runtime comparison with SAT or ILP.

\begin{table}[htbp]
\centering\small
\caption{Construction outcomes on 207 new instances.}
\label{tab:validation-methods}
\begin{tabular}{lrrrr}
\toprule
Method & Feasible & Unknown & Certified & Improved\\\midrule
Greedy & 207 & 0 & 48 & 0\\
Matching, spacing $h$ & 207 & 0 & 22 & 6\\
Random order & 95 & 112 & 22 & 5\\
Insertion order & 191 & 16 & 53 & 57\\
DSATUR, spacing $h$ & 207 & 0 & 58 & 46\\
\bottomrule
\end{tabular}

\tablenote{Certified means that the validated span attains the shared
lower bound. Improved means strictly smaller span than Greedy on the
same instance. Unknown means no order was found within budget; it does
not establish infeasibility. The 207 instances share 69 graphs and are
not 207 independent graph samples.}
\end{table}

\begin{table}[htbp]
\centering\small
\caption{Certificate coverage by graph family.}
\label{tab:validation-families}
\begin{tabular}{lrrrrr}
\toprule
Family & Cases & Greedy & Insertion & DSATUR & Union\\\midrule
BA & 45 & 18 & 7 & 15 & 24\\
Bipartite stress & 9 & 3 & 3 & 3 & 3\\
Cycle stress & 9 & 1 & 0 & 1 & 1\\
ER p=0.15 & 45 & 2 & 0 & 6 & 6\\
ER p=0.30 & 45 & 12 & 43 & 15 & 43\\
Grid stress & 9 & 0 & 0 & 3 & 3\\
Tree & 45 & 12 & 0 & 15 & 16\\
\bottomrule
\end{tabular}

\tablenote{Method columns count certified instances. Union counts each
instance once if any of the five methods attains the shared bound; it
is an aggregate of checked evidence, not a timed portfolio solver.
Each family includes all three parameter pairs. The three stress groups
are deterministic controls and are separate from the random samples.}
\end{table}

Insertion order returns 191 feasible proposals and certifies 53 instances;
random order returns 95 and certifies 22. Excluding the 69 equal-parameter
controls, the ordering success counts are 122/138 and 26/138. The same
compatibility problem is used for $(2,1)$ and $(3,2)$, so these two outcomes
on one graph are dependent. Successful orders of a fixed partition have
the same span $k(r-1)$: the ordering comparison measures the ability to
find an order, not smaller spans among successful orders.

DSATUR-$h$ certifies more instances overall than insertion order (58
versus 53). All 58 DSATUR-$h$ certificates occur for $h=k=1$;
insertion order certifies 22 equal-parameter and 31 unequal-parameter
instances. Spacing all DSATUR colors by $h$ is a conservative feasible
baseline for $h>k$, not the strongest available coloring-based heuristic.
On dense ER, insertion order certifies 43/45 instances and DSATUR-$h$
certifies 15/45. Insertion certifies none of the tree, grid,
cycle or $p=0.15$ ER instances in this study. This pattern is consistent
with the restriction to classes of size at most two: sparse squares can
admit substantially larger color classes. The dense-ER group has one
ordering failure for each unequal parameter pair, and failures also
occur on BA and complete-bipartite stress cases. These observations
support a limited structural role for the construction, not universal
superiority over greedy coloring.

The union of the five methods certifies 96/207 instances. An independent
standard-library audit checks all 907 feasible labelings, all 163 saved
matching certificates, the shared partitions and the complete observation
domain. The remaining 111 instances are uncertified by these constructions
and bounds; neither exact backend was run on this cohort. Data and checker
outputs are archived under \path{results/analysis/certificate_validation/}.
These results are separate from the original 117-instance SAT--ILP
experiment. Its later complete certification in
Appendix~\ref{sec:proof-closure} does not change this new-sample study.
\FloatBarrier
